\section{Kestabilan Lubang Hitam}
\label{sec:stabil}

Berdasarkan seluruh spektrum frekuensi yang diperoleh, baik pada ruang-waktu Schwarzschild maupun Schwarzschild Anti-de Sitter, seluruh mode yang berhasil dihitung memiliki bagian imajiner frekuensi (laju peluruhan) bernilai positif. Hasil tersebut dapat diamati pada Tabel \ref{tab:sch_mu0_n0} hingga Tabel \ref{tab:sads_mu1_n1}, di mana tidak ditemukan satu pun mode dengan $\omega_I<0$. Dengan demikian, perturbasi medan skalar yang diberikan pada kedua geometri tidak mengalami pertumbuhan terhadap waktu, melainkan mengalami peluruhan amplitudo hingga akhirnya kembali menuju keadaan setimbang. Berdasarkan kriteria kestabilan yang telah dijelaskan pada Subbab~\ref{sec:qnm}, hasil tersebut menunjukkan bahwa ruang-waktu Schwarzschild maupun Schwarzschild Anti-de Sitter berada dalam kondisi stabil terhadap perturbasi medan skalar pada seluruh parameter yang ditinjau dalam penelitian ini.

Karakteristik kestabilan pada ruang-waktu Schwarzschild dapat diamati lebih lanjut melalui Gambar \ref{fig:saf_n0} dan Gambar \ref{fig:saf_n1}. Berdasarkan Gambar \ref{fig:saf_n0}(a) dan Gambar \ref{fig:saf_n1}(a), bagian imajiner frekuensi cenderung menurun seiring bertambahnya massa lubang hitam. Hal tersebut menunjukkan bahwa laju peluruhan respon lubang hitam terhadap perturbasi menjadi semakin lambat ketika massa lubang hitam ditingkatkan. Meskipun demikian, seluruh nilai $\omega_I$ tetap berada pada daerah positif sehingga peningkatan massa lubang hitam tidak mengubah sifat kestabilannya. Dengan kata lain, lubang hitam Schwarzschild tetap stabil, namun membutuhkan waktu yang lebih lama untuk meredam perturbasi ketika massanya semakin besar.

Pengaruh massa medan skalar juga terlihat pada Gambar \ref{fig:saf_n0} dan Gambar \ref{fig:saf_n1}, yaitu melalui perbandingan antara kurva medan tak bermassa dan medan bermassa. Pada kasus medan bermassa, bagian imajiner frekuensi cenderung lebih kecil dibandingkan medan tak bermassa sehingga proses peluruhan berlangsung lebih lambat. Selain itu, pada beberapa kombinasi parameter, solusi mode kuasinormal tidak lagi diperoleh ketika massa lubang hitam terus diperbesar. Hasil tersebut menunjukkan bahwa adanya keterbatasan penyelesaian komputasi untuk nilai frekuensi yang dapat diselesaikan terhadap metode yang digunakan.

Berbeda dengan ruang-waktu Schwarzschild, karakteristik peluruhan perturbasi pada ruang-waktu Schwarzschild Anti-de Sitter memperlihatkan kecenderungan yang berlawanan. Berdasarkan Gambar \ref{fig:sads_n0}(a) dan Gambar \ref{fig:sads_n1}(a), bagian imajiner frekuensi meningkat seiring bertambahnya jari-jari horizon peristiwa. Hal ini menunjukkan bahwa amplitudo perturbasi mengalami peluruhan yang semakin cepat pada lubang hitam dengan ukuran horizon yang lebih besar. Dengan demikian, peningkatan jari-jari horizon memperpendek waktu yang diperlukan ruang-waktu Schwarzschild Anti-de Sitter untuk meredam perturbasi dan kembali menuju keadaan setimbang.

Pengaruh massa medan skalar pada ruang-waktu Schwarzschild Anti-de Sitter juga berbeda dibandingkan kasus Schwarzschild. Berdasarkan Gambar \ref{fig:sads_n0} dan Gambar \ref{fig:sads_n1}, penambahan massa medan menyebabkan bagian imajiner frekuensi cenderung meningkat dibandingkan medan tak bermassa. Dengan demikian, perturbasi medan skalar bermassa mengalami peluruhan yang lebih cepat. Berbeda dengan ruang-waktu Schwarzschild, seluruh solusi mode kuasinormal pada ruang-waktu Schwarzschild Anti-de Sitter tetap dapat diperoleh menggunakan metode yang digunakan pada seluruh rentang jari-jari horizon peristiwa yang ditinjau, sehingga tidak ditemukan indikasi hilangnya mode tak bermassa ataupun yang bermassa.

Selain parameter geometri dan massa medan, indeks \textit{overtone} juga memberikan pengaruh terhadap karakteristik peluruhan perturbasi. Berdasarkan hasil pada Tabel \ref{tab:sch_mu0_n0} hingga Tabel \ref{tab:sads_mu1_n1}, \textit{overtone} secara umum memiliki bagian imajiner frekuensi yang lebih besar dibandingkan mode fundamental. Hal tersebut menunjukkan bahwa mode \textit{overtone} mengalami peluruhan lebih cepat sehingga kontribusinya terhadap evolusi perturbasi akan menghilang terlebih dahulu. Oleh karena itu, pada waktu yang cukup lama respons ruang-waktu didominasi oleh mode fundamental yang memiliki laju peluruhan paling lambat.